\clearpage\thispagestyle{empty}
\pdfbookmark[0]{Quellen zur Besprechung von Piazzis Sternverzeichnis}{piazzi1815-sources}
\begin{center}\Large\bfseries Quellen zur Besprechung von Piazzis Sternverzeichnis\end{center}

Die Übertragung folgt Carl Friedrich Gauß’ Besprechung in den \emph{Werken}, Bd.~VI (Göttingen, 1874), S.~590--591. Der Erstabdruck erschien in den \emph{Göttingischen gelehrten Anzeigen}, Stück~91, 10.~Juni 1815, S.~903--904.

Giuseppe Piazzis Katalog erschien 1814 in Palermo in der königlichen Militärdruckerei. Das Exemplar der ETH-Bibliothek, Rar~16115, ist über e-rara zugänglich. In der Präzessionsformel für die Rektaszension steht in der Vorrede ein Plusminuszeichen. Gauß’ Besprechung hat ein Pluszeichen. Der Anhang enthält Joseph Pilatis Reduktionen von Bradleys Beobachtungen für 1756. Die Übertragung folgt den bibliografischen Angaben und Zahlen der Besprechung.

KI-Übertragung aus dem historischen Deutsch in modernes Deutsch. KI-erstellter Satz.

\href{https://gdz.sub.uni-goettingen.de/id/PPN236007467}{Werkausgabe}. \href{https://images.p.sub.uni-goettingen.de/iiif/image/gdz:PPN319721507_1815:00000909/full/1650,/0/default.jpg}{Erstdruck, S.~903}. \href{https://images.p.sub.uni-goettingen.de/iiif/image/gdz:PPN319721507_1815:00000910/full/1650,/0/default.jpg}{S.~904}. \href{https://doi.org/10.3931/e-rara-49143}{Piazzis Katalog}. \href{https://www.e-rara.ch/zut/content/zoom/13476308}{Präzession in der Vorrede}. \href{https://www.e-rara.ch/zut/content/zoom/13476311}{Katalogvergleiche}. \href{https://www.e-rara.ch/zut/content/zoom/13476490}{Pilatis Anhang}. \href{https://images.p.sub.uni-goettingen.de/iiif/image/gdz:PPN319721507_1814:00000190/full/1650,/0/default.jpg}{Gauß’ Anzeige zu Bessel}.

\clearpage
